\section{Fazit und Reflexion}
\label{sec:fazit}

% (1) Übernachtungen ohne Emissionen
%     In der Zielfunktion trägt nur der Transport zum CO2-Term bei. Die
%     Dekodierfunktion für die Hotelwahl (Abschnitt 2.3) kennt deshalb nur
%     w1 und w2, nicht w3. Real verursachen Übernachtungen sehr wohl
%     Emissionen, und zwar je nach Hotelkategorie unterschiedlich hohe.
%     Folge: Das Modell unterschätzt die Klimawirkung einer Reise
%     systematisch und kann keinen Zielkonflikt zwischen Hotelkomfort und
%     CO2 abbilden. Erweiterung wäre ein Emissionswert je Hotel und Nacht,
%     der den Decoder um den w3-Term ergänzt.
%
% (2) Feste Rundreise über alle Ziele
%     Pi ist eine Permutation, also wird jedes Ziel genau einmal besucht und
%     die Reise endet am Startort. Weglassen einzelner Ziele oder Mehrfach-
%     besuche sind nicht darstellbar, könnte aber durchaus relevant sein.
%
% (3) Zeitfenster statt konkreter Abfahrtszeiten
%     Die Diskretisierung in drei Fenster je Verkehrsmittel vereinfacht die
%     Nachbarschaftsstruktur, verwischt aber Unterschiede innerhalb eines
%     Fensters.
%
% (4) Statische Datenbasis
%     Alle Preise und Verbindungen stammen aus Nachschlagetabellen. Reale
%     Verfügbarkeiten und Preise ändern sich während der Planung.

% (5) Grenze der gewichteten Summe gegenueber Pareto
%     Die gewichtete Summe (Skalarisierung) findet je Gewichtung genau eine
%     Loesung, und zwar immer eine auf der konvexen Huelle der erreichbaren
%     Punktmenge. Pareto-optimale Plaene in einer Delle der Front sind fuer
%     KEIN Gewicht w optimal; beim Verschieben von w springt das Ergebnis
%     von einem Randpunkt zum anderen.
%     Gleiche Gewichte helfen NICHT. Gleiche Gewichte heissen "addiere beide Ziele gleich
%     stark", nicht "bevorzuge Ausgewogenheit". Eine Summe bewertet nur
%     den Gesamtwert, nicht die Verteilung auf die Teilziele. 
%
%     ANKNUEPFUNG: Die Abschlussbemerkung stellt die Frage, warum der
%     mittlere Bereich der Paretofront die hoechste Attraktivitaet besitzt.
%     Genau diesen Bereich verfehlt die gewichtete Summe systematisch.
%     QUELLE: Eiben/Smith S. 196 (Skalarisierung und ihre Nachteile) sowie
%     S. 197 (Nichtkonvexitaet und Diskontinuitaeten der Paretofront machen
%     Skalarisierungsverfahren Probleme). 

%     Ausblick: als Studienarbeit könnte man noch so und so ausbauen
